% ===================================================================
\chapter{Introducción}

\section{Qué es PULSO}
PULSO es un software para estudiar el riesgo en redes de pagos digitales. Modela cómo un
problema de liquidez se propaga entre entidades que se pagan entre sí, cómo se
sincronizan sus ciclos de pago y cuándo se producen cascadas de impagos. También aplica
ese modelo al riesgo de crédito de la cartera de un banco: estima, a partir de las
transacciones, si los ciclos de pago de cada cliente están sincronizados con los de
quienes le pagan, y usa esa información para predecir el impago.

El software tiene cuatro formas de uso:

\begin{center}
\begin{tabular}{p{0.26\textwidth}p{0.66\textwidth}}
\toprule
\textbf{Forma de uso} & \textbf{Para qué sirve} \\
\midrule
Aplicación web & Explorar el modelo de forma interactiva, sin programar
(capítulo~\ref{cap:app}). \\
Línea de comandos & Generar los datos, entrenar los modelos y correr el experimento de
crédito (capítulo~\ref{cap:cli}). \\
Biblioteca de Python & Usar las funciones del modelo en programas o notebooks propios
(sección~\ref{sec:biblioteca}). \\
Notebooks & Reproducir las figuras y el análisis de cada tarea
(\ruta{notebooks/}). \\
\bottomrule
\end{tabular}
\end{center}

\begin{aviso}
Todos los datos son sintéticos y los parámetros son ilustrativos, no calibrados con
información real. Los resultados sirven para estudiar mecanismos, no para medir el
riesgo de una institución concreta.
\end{aviso}

\section{A quién está dirigido este manual}
A quien instala, ejecuta o modifica el software. No supone conocimientos previos de
finanzas ni de redes: los conceptos necesarios se explican en la sección~\ref{sec:conceptos}.
Para el fundamento teórico completo, consulte la guía completa
(\ruta{docs/Presentaciones/v1/pulso\_guia\_completa.pdf}).

% ===================================================================
\chapter{Instalación}

\section{Requisitos}
\begin{center}
\begin{tabular}{p{0.24\textwidth}p{0.68\textwidth}}
\toprule
\textbf{Elemento} & \textbf{Requisito} \\
\midrule
Sistema operativo & macOS o Linux. En Windows, use Docker o WSL. \\
Python & 3.12 o superior (lo instala \texttt{uv} si no está). \\
uv & Gestor de entornos y dependencias: \url{https://docs.astral.sh/uv/}. \\
Espacio en disco & Unos 2 GB para el entorno (PyTorch ocupa la mayor parte). \\
Memoria & 8 GB recomendados. \\
Navegador & Cualquier navegador moderno, para la aplicación. \\
Docker (opcional) & Para ejecutar la aplicación sin instalar Python. \\
TeX Live (opcional) & Solo para regenerar los documentos PDF. \\
\bottomrule
\end{tabular}
\end{center}

\section{Instalación con uv}
En una terminal, dentro de la carpeta del proyecto:
\begin{terminal}
uv sync                       # crea el entorno .venv e instala las versiones exactas de uv.lock
uv run pytest -m "not slow"   # comprobación rápida: unas 107 pruebas en ~10 s
\end{terminal}
Si todas las pruebas pasan, la instalación es correcta. La batería completa, con las
pruebas lentas, se corre con \texttt{uv run pytest} y tarda alrededor de un minuto.

\section{Instalación con Docker}
\begin{terminal}
docker build -t pulso .
docker run -p 8501:8501 pulso
\end{terminal}
Abra \texttt{http://localhost:8501}. La imagen incluye los datos y los modelos ya
entrenados, así que la aplicación funciona sin correr ningún experimento.

\section{Primeros pasos}
\begin{enumerate}
\item Abra la aplicación:
\begin{terminal}
uv run streamlit run app/Inicio.py
\end{terminal}
Se abre \texttt{http://localhost:8501}. Para usar otro puerto, agregue
\texttt{--server.port=8502}.
\item En \emph{Inicio}, presione \boton{Cargar escenario de ejemplo “Bre-B estilizado”}.
\item Recorra las páginas en orden con la barra lateral.
\item Para cerrar la aplicación, presione \texttt{Ctrl+C} en la terminal.
\end{enumerate}

\begin{aviso}
En macOS, si el proyecto está dentro de \texttt{\textasciitilde/Documents}, puede aparecer
\texttt{ModuleNotFoundError: No module named 'pulso'}. La solución está en el
capítulo~\ref{cap:problemas}.
\end{aviso}

